% Answer sheet for the Chapter 1 exercises. Build here: pdflatex (run twice).
\documentclass[12pt]{article}
\newcommand{\guideroot}{../}
% Shared preamble for the BARCS student guide (both chapters and the combined book).
% Written to the rules of ASD-STE100 Simplified Technical English:
% short sentences, active voice, present tense, one idea per sentence,
% one meaning per word, and every symbol defined before use.
\usepackage{needspace}
\usepackage[margin=1in]{geometry}
\usepackage[T1]{fontenc}
\usepackage{lmodern}
\usepackage{amsmath,amssymb}
\usepackage{booktabs}
\usepackage{array}
\usepackage{enumitem}
\usepackage[most]{tcolorbox}
\usepackage{xcolor}
\usepackage{tikz}
\usetikzlibrary{arrows.meta,positioning}
\usepackage[hidelinks]{hyperref}

\definecolor{barcsblue}{HTML}{0072B2}
\definecolor{barcsgreen}{HTML}{009E73}
\definecolor{barcsdark}{HTML}{1F3B4D}
\definecolor{barcsorange}{HTML}{D55E00}

\newtcolorbox{symbols}{colback=barcsblue!5,colframe=barcsblue,
  fonttitle=\bfseries,title=What the symbols mean,boxrule=0.6pt,arc=2pt}
\newtcolorbox{worked}[1][]{colback=barcsgreen!5,colframe=barcsgreen!80!black,
  fonttitle=\bfseries,title=Worked example#1,boxrule=0.6pt,arc=2pt}
\newtcolorbox{remember}{colback=barcsorange!5,colframe=barcsorange,
  fonttitle=\bfseries,title=Remember,boxrule=0.6pt,arc=2pt}

\newlist{choices}{enumerate}{1}
\setlist[choices]{label=(\Alph*),leftmargin=2.2em,itemsep=1pt,topsep=2pt}
\setlist{itemsep=2pt}

\setlength{\parskip}{6pt}
\setlength{\parindent}{0pt}
\setcounter{secnumdepth}{0}  % the steps carry their own numbers
\setcounter{tocdepth}{2}
\newif\ifstandalone
\newcommand{\guideauthor}{Prepared for Hyun-Hwan Jeong and Kyu-Won Lee}
\newcommand{\guidedate}{BARCS version 0.2.1 \quad$\cdot$\quad Reading level: grades 7 to 9}
\newcommand{\aboutbox}{%
\begin{center}
\begin{minipage}{0.85\textwidth}
\small
\textbf{About this guide.} This guide explains how the BARCS computer
program works. You need only arithmetic, fractions, percents, powers and
square roots. The guide teaches logarithms when it uses them. Both chapters
use the same small set of numbers, so you can follow every step. The guide
uses the writing rules of ASD-STE100 Simplified Technical English.
\end{minipage}
\end{center}}
% Divider page that opens a chapter in the combined book.
\newcommand{\chapterpart}[3]{%
  \clearpage
  \addcontentsline{toc}{part}{Chapter #1: #2}%
  \vspace*{0.25\textheight}
  \begin{center}
  {\Huge\bfseries Chapter #1\par}\vspace{8pt}
  {\LARGE #2\par}\vspace{6pt}
  {\large #3\par}
  \end{center}
  \clearpage}
\providecommand{\guideroot}{}
% Exercise tags: what kind of thinking a question asks for.
\newcommand{\tagcalc}{\textcolor{barcsblue}{\textsf{\footnotesize[Calculate]}}}
\newcommand{\tagexplain}{\textcolor{barcsgreen!70!black}{\textsf{\footnotesize[Explain]}}}
\newcommand{\tagdecide}{\textcolor{barcsorange}{\textsf{\footnotesize[Decide]}}}
\newcommand{\tagwrite}{\textcolor{barcsdark}{\textsf{\footnotesize[Write]}}}
\newlist{exlist}{enumerate}{1}
\setlist[exlist]{label=\textbf{\arabic*.},leftmargin=2em,itemsep=8pt}
\newlist{parts}{enumerate}{1}
\setlist[parts]{label=(\alph*),leftmargin=2em,itemsep=2pt,topsep=2pt}

\standalonetrue
\title{\textbf{How BARCS Finds Important Genes}\\[6pt]
\Large Answer sheet: Chapter 1 exercises}
\author{\guideauthor}
\date{\guidedate}
\newcommand{\why}[1]{\par\smallskip\noindent\textcolor{barcsgreen!70!black}{\textbf{Why this matters.}} #1}
\begin{document}
\maketitle
\thispagestyle{empty}

Each answer shows the calculation first. ``Why this matters'' explains the
idea behind the question. Small differences in the last digit come from
rounding.

\begin{center}
\small
\begin{tabular}{lcrr}
\toprule
Library & Group & $K$ & $N$\\
\midrule
1 & control & 150 & 1{,}000{,}000\\
2 & control & 210 & 1{,}400{,}000\\
3 & drug    &  60 &   800{,}000\\
4 & drug    &  45 &   600{,}000\\
\bottomrule
\end{tabular}
\end{center}

\section*{Part A: From counts to a change}
\begin{exlist}
\item
\begin{parts}
\item Library 1: $150 \div 1{,}000{,}000 = 0.000150 = 150$ per million.
  Library 2: $210 \div 1{,}400{,}000 = 150$ per million. Library 3:
  $60 \div 800{,}000 = 75$ per million. Library 4: $45 \div 600{,}000 = 75$
  per million.
\item Control: $(150 + 210) \div (1{,}000{,}000 + 1{,}400{,}000) = 360 \div
  2{,}400{,}000 = 150$ per million. Drug: $(60 + 45) \div 1{,}400{,}000 =
  105 \div 1{,}400{,}000 = 75$ per million.
\item Fold change: $75 \div 150 = 0.5$. $\log_2(0.5) = -1$: one halving.
\end{parts}
\why{The proportion removes the size of the library. Then libraries of
different sizes can be compared.}

\item The libraries have different sizes. Library~2 has 1.4 million reads
and library~3 has 0.8 million. A bigger library gives bigger counts for
every guide, even with no biology at all. Per million reads, library~2 has
150 and library~3 has 75. So the drop is $150 \div 75 = 2$ times, not 3.5
times. The extra $3.5 \div 2 = 1.75$ is just the difference in library size
($1.4 \div 0.8 = 1.75$).
\why{Raw counts mix two things: the biology and the depth of reading.
Dividing by $N$ takes the depth out.}

\item
\begin{parts}
\item $0.00015 \div 0.99985 = 0.000150023$. $\ln(0.000150023) = -8.8047$.
\item Reason 1: a proportion must stay between 0 and 1, but a log-odds can be
  any number. So BARCS can add effects ($\beta_0 + \beta_1 + \dots$) without
  ever giving an impossible proportion. Reason 2: on a log scale, ``half''
  and ``double'' are the same size with opposite signs:
  $\ln(0.5) = -0.693$ and $\ln(2) = +0.693$. Changes that multiply become
  changes that add, which is what the design matrix does.
\end{parts}
\why{The design matrix adds up $\beta$ values. That only makes sense on a
scale where effects add and no value is forbidden.}
\end{exlist}

\section*{Part B: Noise}
\begin{exlist}[resume]
\item
\begin{parts}
\item Distances: $-10$, $0$, $+10$, $0$. Sum: $-10 + 0 + 10 + 0 = 0$.
\item Squares: 100, 0, 100, 0. Average: $200 \div 4 = 50$. Standard
  deviation: $\sqrt{50} = 7.07$.
\item The plain distances always add up to 0, because the values above the
  average balance the values below it. So they cannot measure the spread.
  Squares are never negative, so they do not cancel. A second reason:
  squared spreads of independent values add up, and BARCS needs that.
\end{parts}
\why{Every noise formula in BARCS is a variance. This is the reason the
variance uses squares.}

\item
\begin{parts}
\item $1{,}000{,}000 \times 0.00015 \times 0.99985 = 149.98 \approx 150$.
  Standard deviation: $\sqrt{150} = 12.2$ reads.
\item Extra factor: $1 + 999{,}999 \times 0.000001 = 1.999999 \approx 2.0$.
  Total variance: $150 \times 2.0 = 300$. Standard deviation:
  $\sqrt{300} = 17.3$ reads.
\item $1{,}000{,}000 \times 999{,}999 \times 0.000001 \times 0.00015 \times
  0.99985 = 149.98 \approx 150$. Check: $150 + 150 = 300$. Correct.
\end{parts}
\why{Here the biological noise is as large as the counting noise. A model
with counting noise only would make the noise look half as big as it is,
and too many guides would look significant.}

\item
\begin{parts}
\item Two reads from the same library have the correlation 0.000001. If one
  read is guide g7, the next read is very slightly more likely to be g7 too.
  The correlation of one pair is tiny. But a library of a million reads has
  almost a million million pairs, so together they double the variance
  (Exercise~5).
\item $\rho = 0$. If every flask has exactly the same true proportion, there
  is no wobble, and the reads are unrelated.
\end{parts}
\why{$\rho$ is not a free knob. It is a correlation that you can measure,
and it comes from flasks that are truly different.}

\item
\begin{parts}
\item A--B, A--C, A--D, B--A, B--C, B--D, C--A, C--B, C--D, D--A, D--B,
  D--C. That is 12 pairs.
\item $4 \times 3 = 12$. Correct.
\item $N \times N = 16$ would also count A--A, B--B, C--C and D--D. A read
  does not form a pair with itself. Those four ``self'' terms are the
  squares, and they are the counting noise. They are already counted in
  $N \times \mu \times (1-\mu)$.
\end{parts}
\why{The variance of a count is ``$N$ squares $+$ $N \times (N-1)$ pairs''.
The squares give counting noise, the pairs give biological noise.}

\item
\begin{parts}
\item $N = 4$ million: variance $= 599.9 \times (1 + 3{,}999{,}999 \times
  0.000001) = 599.9 \times 5.0 = 2{,}999.5$; SD $= 54.8$ reads; $54.8 \div 4
  = 13.7$ per million.\\
  $N = 2$ million: variance $= 299.96 \times 3.0 = 899.9$; SD $= 30.0$
  reads; $30.0 \div 2 = 15.0$ per million.\\
  $N = 1$ million: variance $= 300$; SD $= 17.3$ reads; 17.3 per million.
\item A: $13.7 \div \sqrt{2} = 9.7$. B: $15.0 \div \sqrt{4} = 7.5$. C:
  $17.3 \div \sqrt{8} = 6.1$ per million.
\item Option C is the most precise. The extra factor grows with $N$ (2, 3
  and 5 for 1, 2 and 4 million reads). Deeper reading shrinks the counting
  noise, but not the biological noise. One library can never be more precise
  than $\sqrt{\rho \times \mu \times (1-\mu)} = 12.2$ per million, however
  deep you read. Only more replicates average the biological noise away.
  More replicates also give more degrees of freedom for the test.
\end{parts}
\why{This is the most practical lesson of Chapter~1: when biological noise
is present, spend money on replicates, not on depth.}
\end{exlist}

\section*{Part C: The design matrix}
\begin{exlist}[resume]
\item
\[
X = \begin{pmatrix} 1 & 0\\ 1 & 0\\ 1 & 0\\ 1 & 1\\ 1 & 1\\ 1 & 1 \end{pmatrix}
\qquad \text{df} = 6 - 2 = 4.
\]
\why{Each extra replicate adds a row but no column, so it adds one degree of
freedom.}

\item
\begin{parts}
\item
\[
X = \begin{pmatrix} 1 & 0 & 0\\ 1 & 0 & 1\\ 1 & 1 & 0\\ 1 & 1 & 1 \end{pmatrix}
\quad \begin{array}{l} \text{donor 1 control}\\ \text{donor 1 drug}\\
\text{donor 2 control}\\ \text{donor 2 drug} \end{array}
\]
\item The ``donor 2'' column takes away the difference between the two
  donors. Then the drug effect is measured inside each donor. Without it,
  the donor difference counts as noise. The SE becomes bigger, and the test
  can miss a real effect. If the donors were not balanced, the donor
  difference could even look like a drug effect.
\item $\text{df} = 4 - 3 = 1$. One degree of freedom gives very little
  information about the noise. The test needs $|t| > 12.71$ for
  $p < 0.05$. More donors would help.
\end{parts}
\why{Every column must have a reason. A column for a known difference
removes noise, but it also costs one degree of freedom.}

\item
\begin{parts}
\item $\eta = 1 \times (-8.8048) + 1 \times (-0.6931) = -9.4979$.
\item $e^{-9.4979} = 0.0000750$. $\mu = 0.0000750 \div 1.0000750 =
  0.0000750 = 75$ per million.
\item $800{,}000 \times 0.0000750 = 60$ reads. The real count is 60. They
  agree.
\end{parts}
\why{A row of $X$ times the $\beta$ values gives a prediction for any
library. This is how BARCS compares the model with the data.}

\item $e^{-0.6931} = 0.50$. The drug halves the odds of guide g7. For small
proportions the odds and the proportion are almost equal, so the proportion
also halves (150 to 75 per million). On the log-odds scale the effect is the
same number for any starting level, and it adds to $\beta_0$ in the design
matrix.
\why{$\beta_1$ is the answer to the biological question: by how much does
the drug change this guide?}
\end{exlist}
\end{document}
